\documentclass[12pt]{article}

\usepackage[T1]{fontenc}
\usepackage[utf8]{inputenc}

\usepackage{amsmath,amssymb,mathtools}
\usepackage{geometry}
\geometry{margin=2.5cm}

\setlength{\parindent}{0pt}
\setlength{\parskip}{4pt}

\newcommand{\dd}{\,\mathrm{d}}
\newcommand{\pd}[2]{\frac{\partial #1}{\partial #2}}
\newcommand{\pdd}[2]{\frac{\partial^2 #1}{\partial #2^2}}

\title{Black--Scholes Derivation}
\author{Egor Shiryaev}
\date{}

\begin{document}
\maketitle

\section{Unhedged Portfolio}

Let the price of the underlying asset $S_t$ be a random process,
and denote the value of the derivative instrument by $V(S,t)$.

Consider a portfolio consisting of a long position in the derivative instrument
and a short position in the underlying asset. At this stage, we treat the
portfolio as unhedged.

Portfolio value:
\begin{equation}
\Pi_t = V(S_t,t) - S_t.
\end{equation}

Change in portfolio value:
\begin{equation}
\dd \Pi_t = \dd V(S_t,t) - \dd S_t.
\end{equation}

\section{Dynamics of the Underlying Asset}

Assume the stock price follows geometric Brownian motion:
\begin{equation}
\dd S_t = \mu S_t \dd t + \sigma S_t \dd W_t,
\end{equation}
where $\mu$ is the drift, $\sigma$ is the volatility,
and $W_t$ is a Wiener process.

Thus the random component introduces uncertainty into our portfolio through the
short position in the underlying asset and the long position in the derivative
instrument.

\section{Dynamics of the Derivative Instrument}

It\^o's lemma lets us find the differential of a function depending on time
and a stochastic process:
\begin{equation}
\dd V = \pd{V}{t}\dd t + \pd{V}{S}\dd S + \frac{1}{2}\pdd{V}{S}(\dd S)^2.
\end{equation}

We take the small change in the value of our stock derivative to be given by
expression (4).

Let us find $(\dd S)^2$ from the price dynamics:
\[
(\dd S)^2 = (\mu S \dd t + \sigma S \dd W)^2.
\]

Expanding:
\[
(\dd S)^2 =
\mu^2 S^2 \dd t^2
+ 2\mu\sigma S^2 \dd t\,\dd W
+ \sigma^2 S^2 (\dd W)^2.
\]

Using the rules of stochastic calculus:
\[
(\dd W)^2=\dd t,\qquad
\dd W\,\dd t=0,\qquad
\dd t^2=0,
\]
we obtain:
\begin{equation}
(\dd S)^2=\sigma^2 S^2 \dd t.
\end{equation}

Substituting:
\begin{equation}
\dd V =
\pd{V}{t}\dd t
+ \pd{V}{S}(\mu S \dd t + \sigma S \dd W)
+ \frac{1}{2}\sigma^2 S^2 \pdd{V}{S}\dd t.
\end{equation}

\section{Source of Randomness}

At first glance, it may seem that two terms in the equation introduce an
element of randomness into the change in the derivative's value:
\[
\frac{1}{2}\frac{\partial^2 V}{\partial S^2}(\dd S)^2,
\qquad
\frac{\partial V}{\partial S}\dd S.
\]

Consider the first term. Using the relation
\[
(\dd S)^2=\sigma^2 S^2 \dd t,
\]
which follows from It\^o's formula (see (4)--(5)), we find that this term is
of order $\dd t$ and is therefore deterministic.

The second term contains the increment of the process $S_t$. Substituting
the underlying asset's dynamics
\[
\dd S = \mu S \dd t + \sigma S \dd W,
\]
we see that this is precisely where the stochastic term arises,
proportional to $\dd W$.

Thus, the sole source of randomness in the dynamics of the derivative's
value is the term
\[
\sigma S \frac{\partial V}{\partial S}\dd W.
\]

\section{Introducing the Hedge}

Since the unhedged portfolio remains random, we can choose the number of
shares so as to eliminate the stochastic term.

Introduce the number of shares $\Delta_t$ and consider the portfolio:
\[
\Pi_t = V(S_t,t) - \Delta_t S_t.
\]

Then:
\[
\dd \Pi_t = \dd V - \Delta_t \dd S_t.
\]

Substituting the expressions:
\[
\dd \Pi =
\Bigl(\pd{V}{t} + \mu S \pd{V}{S} + \frac{1}{2}\sigma^2 S^2 \pdd{V}{S}\Bigr)\dd t
+ \sigma S \pd{V}{S}\dd W
- \Delta_t(\mu S\dd t + \sigma S\dd W).
\]

Group the random terms:
\[
\dd \Pi =
\Bigl(\pd{V}{t} + \mu S \pd{V}{S}
+ \frac{1}{2}\sigma^2 S^2 \pdd{V}{S}
- \Delta_t\mu S\Bigr)\dd t
+ \sigma S(\pd{V}{S}-\Delta_t)\dd W.
\]

The random term vanishes when we choose
\begin{equation}
\Delta_t=\pd{V}{S}.
\end{equation}

This fully eliminates the stochastic risk. The portfolio becomes locally
risk-free:
\begin{equation}
\dd \Pi =
\Bigl(\pd{V}{t} + \frac{1}{2}\sigma^2 S^2 \pdd{V}{S}\Bigr)\dd t.
\end{equation}

Since the portfolio $\Pi_t$ contains no stochastic term, its dynamics
become deterministic. Under no-arbitrage, any risk-free portfolio must
earn the risk-free interest rate $r$; otherwise it would be possible to
earn a guaranteed profit without risk.

Consequently, the dynamics of such a portfolio must satisfy the growth
equation of a risk-free asset.

\section{The Black--Scholes Equation}

A risk-free portfolio must earn the risk-free return $r$:
\begin{equation}
\dd \Pi = r\Pi \dd t.
\end{equation}

Notice that the parameter $\mu$, which characterizes the stock's expected
return, drops out of the subsequent calculations. This means the
derivative's price does not depend on the asset's real expected return.
Equivalently, the calculations are carried out under the risk-neutral
measure $Q$, under which the underlying asset's dynamics take the form
\[
\dd S = r S \dd t + \sigma S \dd W^Q .
\]

Substituting the portfolio:
\[
\dd \Pi = r\bigl(V - S\pd{V}{S}\bigr)\dd t.
\]

Equating:
\[
\pd{V}{t} + \frac{1}{2}\sigma^2 S^2 \pdd{V}{S}
= r\bigl(V - S\pd{V}{S}\bigr),
\]
we obtain:
\begin{equation}
\boxed{
\pd{V}{t}
+ \frac{1}{2}\sigma^2 S^2 \pdd{V}{S}
+ rS\pd{V}{S}
- rV = 0
}
\end{equation}

This is the Black--Scholes equation.

\section{Closed-Form Solution}

So we have obtained the Black--Scholes equation in differential form.
What next?

We often encounter the following formula:
\[
C = S\cdot N(d_1) - K e^{-rT} N(d_2),
\]
\[
d_1 = \frac{\ln\left(\frac{S}{K}\right) + \left(r + \frac12\sigma^2\right)T}{\sigma\sqrt{T}},
\qquad
d_2 = d_1 - \sigma\sqrt{T}.
\]

This is the Black--Scholes formula for a European call option
(no dividends).

What does it mean: that's the first question worth asking when looking
at this formula.

We know that for the option contract to be exercised, the condition
$(S_T - K)^+$ must hold, in other words
\begin{equation}
S_T - K > 0 \Rightarrow S_T > K.
\end{equation}

At this point it is worth taking the logarithm of the expression.
We get
\[
\ln(S_T) - \ln(K).
\]

Now we need to go back and recall that we assumed pricing follows
geometric Brownian motion.

Under the real (physical) measure $P$, the asset's dynamics are
\[
\dd S = \mu S \dd t + \sigma S \dd W.
\]

However, when pricing derivatives we work under the risk-neutral
measure $Q$, in which the asset's expected return equals the
risk-free rate $r$. Under this measure the dynamics become
\[
\dd S = r S \dd t + \sigma S \dd W^Q.
\]

Integrating this equation, we get
\[
\ln(S_T) = \ln(S_0) + \left(r - \frac12\sigma^2\right)T + \sigma W_T.
\]

This gives:
\[
\ln(S_0) + \left(r - \frac12\sigma^2\right)T + \sigma W_T > \ln(K),
\]
then
\[
\ln\left(\frac{K}{S_0}\right)
<
\left(r - \frac12\sigma^2\right)T + \sigma W_T.
\]

Since $W_T\sim N(0,T)$, we standardize:
\[
W_T=\sqrt{T}\,Z, \qquad Z\sim N(0,1).
\]

Substituting:
\[
\ln\left(\frac{K}{S_0}\right)
<
\left(r - \frac12\sigma^2\right)T + \sigma\sqrt{T}Z.
\]

Then:
\[
\sigma\sqrt{T}Z >
\ln\left(\frac{K}{S_0}\right)
-
\left(r - \frac12\sigma^2\right)T,
\]
\[
Z >
\frac{1}{\sigma\sqrt{T}}
\left(
\ln\left(\frac{K}{S_0}\right)
-
\left(r - \frac12\sigma^2\right)T
\right).
\]

Let us pause here. Earlier we said that $Z\sim N(0,1)$.

To understand this we must answer: what is $N(x)$? $N(x)=P(Z\le x)$.

This is the area under the normal distribution curve to the left of
the point $x$. In our case we want the call option to be exercised,
so we're interested in the area to the right of the point.

Denote this point by $d_2$; then
\[
P(Z>-d_2).
\]

Flipping the expression:
\[
Z >
\frac{1}{\sigma\sqrt{T}}
\left(
\ln\left(\frac{S_0}{K}\right)
+
\left(r - \frac12\sigma^2\right)T
\right)
= -d_2.
\]

In other words:
\begin{equation}
P(S_T>K)=N(d_2).
\end{equation}

Recall that we are working under the risk-neutral measure, so the
option price is:
\begin{equation}
C_0=e^{-rT}\mathbb{E}^Q[(S_T-K)^+].
\end{equation}

Let's break this down further:
\begin{equation}
C_0=e^{-rT}\mathbb{E}^Q[S_T\mathbf{1}_{S_T>K}-K\mathbf{1}_{S_T>K}],
\end{equation}

\begin{equation}
C_0=e^{-rT}
\left(
\mathbb{E}^Q[S_T\mathbf{1}_{S_T>K}]
-
K P(S_T>K)
\right).
\end{equation}

We have already introduced the second part and know how to work
with it.

\[
\mathbb{E}^Q[S_T\mathbf{1}_{S_T>K}],\qquad
\text{and we already know how } S_T \text{ behaves:}
\]
\[
S_T=S_0 e^{(r-\frac12\sigma^2)T+\sigma\sqrt{T}Z}.
\]

This immediately gives the expectation:
\[
S_0 e^{(r-\frac12\sigma^2)T}
\mathbb{E}^Q\left[e^{\sigma\sqrt{T}Z}\mathbf{1}_{Z>-d_2}\right].
\]

From this we get the integral:
\[
\int_{-d_2}^{\infty}
e^{\sigma\sqrt{T}z}
\frac{1}{\sqrt{2\pi}}e^{-z^2/2}\dd z.
\]

For convenience, denote:
\[
\theta=\sigma\sqrt{T}.
\]

\[
\int_{-d_2}^{\infty}
\frac{1}{\sqrt{2\pi}}
e^{\theta z-\frac12 z^2}\dd z.
\]

Completing the square:
\[
e^{\frac12\theta^2}
\int_{-d_2-\theta}^{\infty}
\frac{1}{\sqrt{2\pi}}e^{-q^2/2}\dd q.
\]

Substitution:
\[
q=z-\theta.
\]

Note that
\[
d_1=d_2+\sigma\sqrt{T}=d_2+\theta.
\]

Hence the lower limit of integration transforms as:
\[
-d_2-\theta=-d_1.
\]

Then
\[
\int_{-d_2-\theta}^{\infty}\phi(q)\dd q
=
P(Z>-d_1)=N(d_1).
\]

Finally:
\[
\mathbb{E}^Q[S_T\mathbf{1}_{S_T>K}]
=
S_0 e^{rT}N(d_1).
\]

Returning to our expression:
\begin{equation}
\fbox{$\displaystyle
\begin{aligned}
C_0 &= S_0N(d_1) - K e^{-rT}N(d_2),\\
\text{where}\quad d_1 &= \frac{\ln\left(\frac{S_0}{K}\right) + \left(r+\frac12\sigma^2\right)T}{\sigma\sqrt{T}},\\
d_2 &= d_1 - \sigma\sqrt{T}.
\end{aligned}
$}
\end{equation}
\end{document}
